% XeLaTeX can use any Mac OS X font. See the setromanfont command below.
% Input to XeLaTeX is full Unicode, so Unicode characters can be typed directly into the source.
% The next lines tell TeXShop to typeset with xelatex, and to open and save the source with Unicode encoding.
%!TEX TS-program = xelatex
%!TEX encoding = UTF-8 Unicode

\documentclass[12pt]{article}
\usepackage[landscape, margin=.4in]{geometry}              % See geometry.pdf to learn the layout options. There are lots.
\usepackage[parfill]{parskip}    % Activate to begin paragraphs with an empty line rather than an indent
\usepackage{graphicx}
\usepackage{ifthen}
\usepackage{titlesec}
\usepackage{xcolor}
\usepackage{array}
\usepackage{xstring}

% Will Robertson's fontspec.sty can be used to simplify font choices.
% To experiment, open /Applications/Font Book to examine the fonts provided on Mac OS X,
% and change "Hoefler Text" to any of these choices.
\usepackage{fontspec,xltxtra,xunicode}
\defaultfontfeatures{Mapping=tex-text}
\setromanfont[Mapping=tex-text]{Nanum Myeongjo}
\setsansfont[Scale=MatchLowercase,Mapping=tex-text]{Euphemia UCAS}
\setmonofont[Scale=MatchLowercase]{SimHei}

\pagestyle{empty}
\setlength{\parindent}{0pt}

\newcommand{\messagename}{Message Name}
\newcommand{\cipher}{Caesar }
\newcommand{\key}{3}
%numbers is true if the message comes in numbers
\newboolean{numbers}
\setboolean{numbers}{false}
%numtable is true if you want numbers in the ciphertable
\newboolean{numtable}
\setboolean{numtable}{true}
%numpages is the number of pages needed for the message
\newcount\numpages
\numpages=2
\newcommand{\ccmessagefontsize}{14}%fontsize for message letters
\newcommand{\ccmessagebreak}{18pt}%space between lines in the message
\newcommand{\ccmessagenumfontsize}{12}%fontsize for message letters
\newcommand{\cctablefontsize}{14}%fontsize for cipher table letters
\newcommand{\nfs}{\fontsize{9}{9}\selectfont}%fontsize for cipher table numbers
\newcommand{\lfs}{\fontsize{10}{10}\selectfont\texttt}%fontsize for cipher table numbers
\newcommand{\topmessageskip}{\vskip 28pt}
\newcounter{chrnum}
\newcounter{chrnumtwo}
%strlen is the number of chacters that prints in one line of the message
\newcount \strlen 
\ifthenelse{\boolean{numbers}}{\strlen=99}{\strlen=40}
%    \numpages=4
\newcommand{\oneone}{abcde fghij klmno pqrst uvwxy z}
\newcommand{\onetwo}{00 01 02 03 04 \ \ \ 05 06 07 08 09 \ \ \ 10 11 12 13 14 \ \ \ 15 16 17 18 19 \ \ \  20 21 22 23 24 \ \ \ 25}
\newcommand{\twoone}{abcde fghij klmno pqrst uvwxy z}
\newcommand{\twotwo}{00 01 02 03 04 \ \ \ 05 06 07 08 09 \ \ \ 10 11 12 13 14 \ \ \ 15 16 17 18 19 \ \ \  20 21 22 23 24 \ \ \ 25}
\newcommand{\threeone}{ } 
\newcommand{\threetwo}{JFDVNYVIV ZE SVUWFIU TFLEKP, MZIXZEZR.}
\newcommand{\fourone}{ }
\newcommand{\fourtwo}{YV CVWK VETIPGKVU DVJJRXVJ KYRK IVMVRC KYV }
\newcommand{\fiveone}{ }
\newcommand{\fivetwo}{CFTRKZFE. KYVP YRMV EVMVI SVVE TIRTBVU. XF }
\newcommand{\sixone}{ }
\newcommand{\sixtwo}{KF KYV TFDZTJ JVTKZFE FW TIPGKFTCLS.FIX KF}
\newcommand{\sevenone}{ }
\newcommand{\seventwo}{IVRU KYV NYFCV JKFIP.}
\newcommand{\eightone}{ }
\newcommand{\eighttwo}{KNF YLEUIVU PVRIJ RXF, KYFDRJ A. SVRCV }
\newcommand{\nineone}{ }
\newcommand{\ninetwo}{SLIZVU R KIVRJLIV FW XFCU REU JZCMVI }
\newcommand{\tenone}{ } 
\newcommand{\tentwo}{JFDVNYVIV ZE SVUWFIU TFLEKP, MZIXZEZR.}
\newcommand{\elevenone}{ }
\newcommand{\eleventwo}{YV CVWK VETIPGKVU DVJJRXVJ KYRK IVMVRC KYV }
\newcommand{\twelveone}{ }
\newcommand{\twelvetwo}{CFTRKZFE. KYVP YRMV EVMVI SVVE TIRTBVU. XF }
\newcommand{\thirteenone}{ }
\newcommand{\thirteentwo}{KF KYV TFDZTJ JVTKZFE FW TIPGKFTCLS.FIX KF}
\newcommand{\fourteenone}{ }
\newcommand{\fourteentwo}{IVRU KYV NYFCV JKFIP.}
%
\newcommand{\stripskip}{-1pt}
\renewcommand{\arraystretch}{1.2}
\newcommand{\redcell}[1]{\textcolor{red}{\hbox to 15pt{#1\hfil}}} 
\newcommand{\redcellnum}[1]{\fontsize{12}{12}\textcolor{red}{\hbox to 1.5em{#1\hfil}}} 
\newcommand{\cell}[1]{\hbox to 15pt{#1\hfil}}
\newcommand{\tcell}[1]{\nfs{\hbox to 5pt{#1\hfil}}}

\begin{document}
%top banner in three horizontal boxes
\begin{minipage}[b][.75in][t]{3in}{\includegraphics[scale = 2]{logo_small}}\end{minipage}%
\begin{minipage}[b][.75in][c]{5in}{\vskip10pt\fontsize{24}{26}\selectfont{Crack Substitution}}\end{minipage}%
\begin{minipage}[b][.75in][c]{2in}{{\bf{\messagename}}\\\cipher{ Cipher}}\end{minipage}%
\hrule

%Left side of page has three boxes: ciphertable, message, instructions
%\fbox{
\begin{minipage}[b][6.55in][t]{8.5in}
%%ciphertable: has two rows for letters and four for numtable
%    \fbox{
    \begin{minipage}[b][1.25in][c]{8.3in}
    {
    \renewcommand{\arraystretch}{1.5}%sets stretch in height of row
    \ifthenelse{\boolean{numtable}}{\setlength\tabcolsep{5.5pt}}{}%increases space between columns
    \fontsize{\cctablefontsize}{\cctablefontsize}\selectfont\texttt{%
    \begin{tabular}{r |*{26}{c|}}
    \cline{2 - 27}%
    {\fontsize{10}{10}\selectfont{plaintext}}
&a&b&c&d&e&f&g&h&i&j&k&l&m&n&o&p&q&r&s&t&u&v&w&x&y&z\\
    \cline{2 - 27}% 
\ifthenelse{\boolean{numtable}}{
{\fontsize{10}{10}\selectfont{numbers}}&\nfs{ 0}&\nfs{1}&\nfs{2}&\nfs{3}&\nfs{4}&\nfs{5}&\nfs{6}&\nfs{7}&\nfs{8}&\nfs{9}&\nfs{10}&\nfs{11}&\nfs{12}&\nfs{13}&\nfs{14}&\nfs{15}&\nfs{16}&\nfs{17}&\nfs{18}&\nfs{19}&\nfs{20}&\nfs{21}&\nfs{22}&\nfs{23}&\nfs{24}&\nfs{25}\\%
\cline{2 - 27}%
{\fontsize{10}{10}\selectfont{CIPHERNUMBERS}}&&&&&&&&&&&&&&&&&&&&&&&&&&\\
\cline{2 - 27}%
}{}
{\fontsize{10}{10}\selectfont{CIPHERTEXT}}&&&&&&&&&&&&&&&&&&&&&&&&&&\\
 \cline{2 - 27}
    \end{tabular}%
    }}\end{minipage}
%    }		
\ifthenelse{\boolean{numbers}}{%
\begin{minipage}[b][4.8in][t]{8.3in}%
{\topmessageskip%
\fontsize{\ccmessagenumfontsize}{\ccmessagenumfontsize}\selectfont\texttt{%
\setcounter{chrnum}{1}\setcounter{chrnumtwo}{2}%
\loop\redcellnum{\StrChar{\oneone}{\thechrnum}}\ifnum\thechrnum<\strlen\stepcounter{chrnum}\repeat\\%
\setcounter{chrnum}{1}\setcounter{chrnumtwo}{2}%%
\loop\StrMid{\onetwo}{\thechrnum}{\thechrnumtwo}\ \ifnum\thechrnum<\strlen\addtocounter{chrnum}{3}\addtocounter{chrnumtwo}{3}\repeat\\[\ccmessagebreak]%
\setcounter{chrnum}{1}%
\loop\redcellnum{\StrChar{\twoone}{\thechrnum}}\ifnum\thechrnum<\strlen\stepcounter{chrnum}\repeat\\%
\setcounter{chrnum}{1}\setcounter{chrnumtwo}{2}%%
\loop\StrMid{\twotwo}{\thechrnum}{\thechrnumtwo}\ \ifnum\thechrnum<\strlen\addtocounter{chrnum}{3}\addtocounter{chrnumtwo}{3}\repeat\\[\ccmessagebreak]%
\setcounter{chrnum}{1}%
\loop\redcellnum{\StrChar{\threeone}{\thechrnum}}\ifnum\thechrnum<\strlen\stepcounter{chrnum}\repeat\\%
\setcounter{chrnum}{1}\setcounter{chrnumtwo}{2}%%
\loop\StrMid{\threetwo}{\thechrnum}{\thechrnumtwo}\ \ifnum\thechrnum<\strlen\addtocounter{chrnum}{3}\addtocounter{chrnumtwo}{3}\repeat\\[\ccmessagebreak]%
\setcounter{chrnum}{1}%
\loop\redcellnum{\StrChar{\fourone}{\thechrnum}}\ifnum\thechrnum<\strlen\stepcounter{chrnum}\repeat\\%
\setcounter{chrnum}{1}\setcounter{chrnumtwo}{2}%%
\loop\StrMid{\fourtwo}{\thechrnum}{\thechrnumtwo}\ \ifnum\thechrnum<\strlen\addtocounter{chrnum}{3}\addtocounter{chrnumtwo}{3}\repeat\\[\ccmessagebreak]%
\setcounter{chrnum}{1}%
\loop\redcellnum{\StrChar{\fiveone}{\thechrnum}}\ifnum\thechrnum<\strlen\stepcounter{chrnum}\repeat\\%
\setcounter{chrnum}{1}\setcounter{chrnumtwo}{2}%%
\loop\StrMid{\fivetwo}{\thechrnum}{\thechrnumtwo}\ \ifnum\thechrnum<\strlen\addtocounter{chrnum}{3}\addtocounter{chrnumtwo}{3}\repeat\\[\ccmessagebreak]%
\loop\redcellnum{\StrChar{\sevenone}{\thechrnum}}\ifnum\thechrnum<\strlen\stepcounter{chrnum}\repeat\\%
\setcounter{chrnum}{1}\setcounter{chrnumtwo}{2}%
\loop\StrMid{\seventwo}{\thechrnum}{\thechrnumtwo}\ \ifnum\thechrnum<\strlen\addtocounter{chrnum}{3}\addtocounter{chrnumtwo}{3}\repeat}%
}\end{minipage}
%    }
}{
\begin{minipage}[b][4.8in][t]{8.3in}%
{\topmessageskip%
\fontsize{\ccmessagefontsize}{\ccmessagefontsize}\selectfont\texttt{%
\setcounter{chrnum}{1}%
\loop\redcell{\StrChar{\oneone}{\thechrnum}}\ifnum\thechrnum<\strlen\stepcounter{chrnum}\repeat\\%
\setcounter{chrnum}{1}%
\loop\cell{\StrChar{\onetwo}{\thechrnum}}\ifnum\thechrnum<\strlen\stepcounter{chrnum}\repeat\\[\ccmessagebreak]%
\setcounter{chrnum}{1}%
\loop\redcell{\StrChar{\twoone}{\thechrnum}}\ifnum\thechrnum<\strlen\stepcounter{chrnum}\repeat\\%
\setcounter{chrnum}{1}%
\loop\cell{\StrChar{\twotwo}{\thechrnum}}\ifnum\thechrnum<\strlen\stepcounter{chrnum}\repeat\\[\ccmessagebreak]%
\setcounter{chrnum}{1}%
\loop\redcell{\StrChar{\threeone}{\thechrnum}}\ifnum\thechrnum<\strlen\stepcounter{chrnum}\repeat\\%
\setcounter{chrnum}{1}%
\loop\cell{\StrChar{\threetwo}{\thechrnum}}\ifnum\thechrnum<\strlen\stepcounter{chrnum}\repeat\\[\ccmessagebreak]%
\setcounter{chrnum}{1}%
\loop\redcell{\StrChar{\fourone}{\thechrnum}}\ifnum\thechrnum<\strlen\stepcounter{chrnum}\repeat\\%
\setcounter{chrnum}{1}%
\loop\cell{\StrChar{\fourtwo}{\thechrnum}}\ifnum\thechrnum<\strlen\stepcounter{chrnum}\repeat\\[\ccmessagebreak]%
\setcounter{chrnum}{1}%
\loop\redcell{\StrChar{\fiveone}{\thechrnum}}\ifnum\thechrnum<\strlen\stepcounter{chrnum}\repeat\\%
\setcounter{chrnum}{1}%
\loop\cell{\StrChar{\fivetwo}{\thechrnum}}\ifnum\thechrnum<\strlen\stepcounter{chrnum}\repeat\\[\ccmessagebreak]%
\setcounter{chrnum}{1}%
\loop\redcell{\StrChar{\sixone}{\thechrnum}}\ifnum\thechrnum<\strlen\stepcounter{chrnum}\repeat\\%
\setcounter{chrnum}{1}%
\loop\cell{\StrChar{\sixtwo}{\thechrnum}}\ifnum\thechrnum<\strlen\stepcounter{chrnum}\repeat\\[\ccmessagebreak]%
\setcounter{chrnum}{1}%
\loop\redcell{\StrChar{\sevenone}{\thechrnum}}\ifnum\thechrnum<\strlen\stepcounter{chrnum}\repeat\\%
\setcounter{chrnum}{1}%
\loop\cell{\StrChar{\seventwo}{\thechrnum}}\ifnum\thechrnum<\strlen\stepcounter{chrnum}\repeat}%
}\end{minipage}
%    }
}
%Instruction box
\fbox{%
\begin{minipage}[b][.6in][t]{8.3in}%
{\fontsize{11}{11}\selectfont\bf{To Crack a Substitution Cipher:}}%
{\fontsize{11}{11}\selectfont{ Enter substitutions in the table and above the message.%
\\Use letter frequencies to help. Use patterns in English  or patterns in the cipher table to help.% 
\\Make changes until your message makes sense.}}%
\end{minipage}%
}
\end{minipage}
%}%
%\fbox{
\begin{minipage}[b][6.55in][t]{1.7in}
\renewcommand{\arraystretch}{1.5}%sets stretch in height of row
\setlength\tabcolsep{1pt}%increases space between columns
\vskip 20pt
\hfil\lfs{Letter Frequencies}\hfil
{\center{\begin{tabular}{| c | c | c | c |}
	\hline
    	\lfs{In message}&\lfs{In English}\\
    	\hline
    	\lfs V -\nfs{15.4}&  \lfs e -\nfs{12.7}\\
    	\lfs F - \nfs{9.3}&  \lfs t -\  \nfs{9.1}\\
    	\lfs K - \nfs{8.4}&  \lfs a - \nfs{8.2}\\
    	\lfs I - \nfs{7.9}&  \lfs o - \nfs{7.5}\\
    	\lfs R - \nfs{7.0}&  \lfs i - \nfs{7.0}\\
    	\lfs J - \nfs{5.1}&  \lfs n - \nfs{6.8}\\
    	\lfs Y - \nfs{5.1}&  \lfs s - \nfs{6.3}\\
    	\lfs E - \nfs{4.7}&  \lfs h - \nfs{6.1}\\
    	\lfs T - \nfs{4.7}&  \lfs r - \nfs{6.0}\\
    	\lfs U - \nfs{4.7}&  \lfs d - \nfs{4.3}\\
    	\lfs Z - \nfs{4.2}&  \lfs l - \nfs{4.0}\\
    	\lfs C - \nfs{3.7}&  \lfs u - \nfs{2.8}\\
    	\lfs P - \nfs{2.8}&  \lfs c - \nfs{2.8}\\
    	\lfs X - \nfs{2.8}&  \lfs w - \nfs{2.4}\\
    	\lfs L - \nfs{2.3}&  \lfs m - \nfs{2.4}\\
    	\lfs M - \nfs{2.3}&  \lfs f - \nfs{2.2}\\
    	\lfs S - \nfs{2.3}&  \lfs y - \nfs{2.0}\\
    	\lfs D - \nfs{1.9}&  \lfs g - \nfs{2.0}\\
    	\lfs W - \nfs{1.9}&  \lfs p - \nfs{1.9}\\
    	\lfs N - \nfs{1.4}&  \lfs b - \nfs{1.5}\\
    	\lfs G - \nfs{0.9}&  \lfs v - \nfs{1.0}\\
    	\lfs A - \nfs{0.5}&  \lfs k - \nfs{0.8}\\
    	\lfs B - \nfs{0.0}&  \lfs x - \nfs{0.2}\\
    	\lfs H - \nfs{0.0}&  \lfs j - \nfs{0.2}\\
    	\lfs O - \nfs{0.0}&  \lfs z - \nfs{0.1}\\
    	\lfs Q - \nfs{0.0}&  \lfs q - \nfs{0.1}\\
    	\hline
\end{tabular}}\\
}
\end{minipage}
%}
\ifnum\numpages=1{\end{document}}\else{\newpage
\renewcommand{\oneone}{abcde fghij klmno pqrst uvwxy z}
\renewcommand{\onetwo}{00 01 02 03 04 \ \ \ 05 06 07 08 09 \ \ \ 10 11 12 13 14 \ \ \ 15 16 17 18 19 \ \ \  20 21 22 23 24 \ \ \ 25}
\renewcommand{\twoone}{abcde fghij klmno pqrst uvwxy z}
\renewcommand{\twotwo}{00 01 02 03 04 \ \ \ 05 06 07 08 09 \ \ \ 10 11 12 13 14 \ \ \ 15 16 17 18 19 \ \ \  20 21 22 23 24 \ \ \ 25}
\renewcommand{\threeone}{ } 
\renewcommand{\threetwo}{JFDVNYVIV ZE SVUWFIU TFLEKP, MZIXZEZR.}
\renewcommand{\fourone}{ }
\renewcommand{\fourtwo}{YV CVWK VETIPGKVU DVJJRXVJ KYRK IVMVRC KYV }
\renewcommand{\fiveone}{ }
\renewcommand{\fivetwo}{CFTRKZFE. KYVP YRMV EVMVI SVVE TIRTBVU. XF }
\renewcommand{\sixone}{ }
\renewcommand{\sixtwo}{KF KYV TFDZTJ JVTKZFE FW TIPGKFTCLS.FIX KF}
\renewcommand{\sevenone}{ }
\renewcommand{\seventwo}{IVRU KYV NYFCV JKFIP.}
\renewcommand{\eightone}{ }
\renewcommand{\eighttwo}{KNF YLEUIVU PVRIJ RXF, KYFDRJ A. SVRCV }
\renewcommand{\nineone}{ }
\renewcommand{\ninetwo}{SLIZVU R KIVRJLIV FW XFCU REU JZCMVI }
\renewcommand{\tenone}{ } 
\renewcommand{\tentwo}{JFDVNYVIV ZE SVUWFIU TFLEKP, MZIXZEZR.}
\renewcommand{\elevenone}{ }
\renewcommand{\eleventwo}{YV CVWK VETIPGKVU DVJJRXVJ KYRK IVMVRC KYV }
\renewcommand{\twelveone}{ }
\renewcommand{\twelvetwo}{CFTRKZFE. KYVP YRMV EVMVI SVVE TIRTBVU. XF }
\renewcommand{\thirteenone}{ }
\renewcommand{\thirteentwo}{KF KYV TFDZTJ JVTKZFE FW TIPGKFTCLS.FIX KF}
\renewcommand{\fourteenone}{ }
\renewcommand{\fourteentwo}{IVRU KYV NYFCV JKFIP.}
\ifthenelse{\boolean{numbers}}{
\begin{minipage}[b][4.8in][t]{8.3in}%
{\topmessageskip%
\fontsize{\ccmessagenumfontsize}{\ccmessagenumfontsize}\selectfont\texttt{%
\setcounter{chrnum}{1}\setcounter{chrnumtwo}{2}%
\loop\redcellnum{\StrChar{\oneone}{\thechrnum}}\ifnum\thechrnum<\strlen\stepcounter{chrnum}\repeat\\%
\setcounter{chrnum}{1}\setcounter{chrnumtwo}{2}%%
\loop\StrMid{\onetwo}{\thechrnum}{\thechrnumtwo}\ \ifnum\thechrnum<\strlen\addtocounter{chrnum}{3}\addtocounter{chrnumtwo}{3}\repeat\\[\ccmessagebreak]%
\setcounter{chrnum}{1}%
\loop\redcellnum{\StrChar{\twoone}{\thechrnum}}\ifnum\thechrnum<\strlen\stepcounter{chrnum}\repeat\\%
\setcounter{chrnum}{1}\setcounter{chrnumtwo}{2}%%
\loop\StrMid{\twotwo}{\thechrnum}{\thechrnumtwo}\ \ifnum\thechrnum<\strlen\addtocounter{chrnum}{3}\addtocounter{chrnumtwo}{3}\repeat\\[\ccmessagebreak]%
\setcounter{chrnum}{1}%
\loop\redcellnum{\StrChar{\threeone}{\thechrnum}}\ifnum\thechrnum<\strlen\stepcounter{chrnum}\repeat\\%
\setcounter{chrnum}{1}\setcounter{chrnumtwo}{2}%%
\loop\StrMid{\threetwo}{\thechrnum}{\thechrnumtwo}\ \ifnum\thechrnum<\strlen\addtocounter{chrnum}{3}\addtocounter{chrnumtwo}{3}\repeat\\[\ccmessagebreak]%
\setcounter{chrnum}{1}%
\loop\redcellnum{\StrChar{\fourone}{\thechrnum}}\ifnum\thechrnum<\strlen\stepcounter{chrnum}\repeat\\%
\setcounter{chrnum}{1}\setcounter{chrnumtwo}{2}%%
\loop\StrMid{\fourtwo}{\thechrnum}{\thechrnumtwo}\ \ifnum\thechrnum<\strlen\addtocounter{chrnum}{3}\addtocounter{chrnumtwo}{3}\repeat\\[\ccmessagebreak]%
\setcounter{chrnum}{1}%
\loop\redcellnum{\StrChar{\fiveone}{\thechrnum}}\ifnum\thechrnum<\strlen\stepcounter{chrnum}\repeat\\%
\setcounter{chrnum}{1}\setcounter{chrnumtwo}{2}%%
\loop\StrMid{\fivetwo}{\thechrnum}{\thechrnumtwo}\ \ifnum\thechrnum<\strlen\addtocounter{chrnum}{3}\addtocounter{chrnumtwo}{3}\repeat\\[\ccmessagebreak]%
\loop\redcellnum{\StrChar{\sevenone}{\thechrnum}}\ifnum\thechrnum<\strlen\stepcounter{chrnum}\repeat\\%
\setcounter{chrnum}{1}\setcounter{chrnumtwo}{2}%
\loop\StrMid{\seventwo}{\thechrnum}{\thechrnumtwo}\ \ifnum\thechrnum<\strlen\addtocounter{chrnum}{3}\addtocounter{chrnumtwo}{3}\repeat}%
}\end{minipage}
%    }
}{
\begin{minipage}[b][4.8in][t]{8.3in}%
{\topmessageskip%
\fontsize{\ccmessagefontsize}{\ccmessagefontsize}\selectfont\texttt{%
\setcounter{chrnum}{1}%
\loop\redcell{\StrChar{\oneone}{\thechrnum}}\ifnum\thechrnum<\strlen\stepcounter{chrnum}\repeat\\%
\setcounter{chrnum}{1}%
\loop\cell{\StrChar{\onetwo}{\thechrnum}}\ifnum\thechrnum<\strlen\stepcounter{chrnum}\repeat\\[\ccmessagebreak]%
\setcounter{chrnum}{1}%
\loop\redcell{\StrChar{\twoone}{\thechrnum}}\ifnum\thechrnum<\strlen\stepcounter{chrnum}\repeat\\%
\setcounter{chrnum}{1}%
\loop\cell{\StrChar{\twotwo}{\thechrnum}}\ifnum\thechrnum<\strlen\stepcounter{chrnum}\repeat\\[\ccmessagebreak]%
\setcounter{chrnum}{1}%
\loop\redcell{\StrChar{\threeone}{\thechrnum}}\ifnum\thechrnum<\strlen\stepcounter{chrnum}\repeat\\%
\setcounter{chrnum}{1}%
\loop\cell{\StrChar{\threetwo}{\thechrnum}}\ifnum\thechrnum<\strlen\stepcounter{chrnum}\repeat\\[\ccmessagebreak]%
\setcounter{chrnum}{1}%
\loop\redcell{\StrChar{\fourone}{\thechrnum}}\ifnum\thechrnum<\strlen\stepcounter{chrnum}\repeat\\%
\setcounter{chrnum}{1}%
\loop\cell{\StrChar{\fourtwo}{\thechrnum}}\ifnum\thechrnum<\strlen\stepcounter{chrnum}\repeat\\[\ccmessagebreak]%
\setcounter{chrnum}{1}%
\loop\redcell{\StrChar{\fiveone}{\thechrnum}}\ifnum\thechrnum<\strlen\stepcounter{chrnum}\repeat\\%
\setcounter{chrnum}{1}%
\loop\cell{\StrChar{\fivetwo}{\thechrnum}}\ifnum\thechrnum<\strlen\stepcounter{chrnum}\repeat\\[\ccmessagebreak]%
\setcounter{chrnum}{1}%
\loop\redcell{\StrChar{\sixone}{\thechrnum}}\ifnum\thechrnum<\strlen\stepcounter{chrnum}\repeat\\%
\setcounter{chrnum}{1}%
\loop\cell{\StrChar{\sixtwo}{\thechrnum}}\ifnum\thechrnum<\strlen\stepcounter{chrnum}\repeat\\[\ccmessagebreak]%
\setcounter{chrnum}{1}%
\loop\redcell{\StrChar{\sevenone}{\thechrnum}}\ifnum\thechrnum<\strlen\stepcounter{chrnum}\repeat\\%
\setcounter{chrnum}{1}%
\loop\cell{\StrChar{\seventwo}{\thechrnum}}\ifnum\thechrnum<\strlen\stepcounter{chrnum}\repeat}%
}\end{minipage}
%    }
}
}\fi
\end{document}
%  {\footnotesize{\it{This is the footer}}}
% For many users, the previous commands will be enough.
% If you want to directly input Unicode, add an Input Menu or Keyboard to the menu bar 
% using the International Panel in System Preferences.
% Unicode must be typeset using a font containing the appropriate characters.
% Remove the comment signs below for examples.

% \newfontfamily{\A}{Geeza Pro}
% \newfontfamily{\H}[Scale=0.9]{Lucida Grande}
% \newfontfamily{\J}[Scale=0.85]{Osaka}

% Here are some multilingual Unicode fonts: this is Arabic text: {\A السلام عليكم}, this is Hebrew: {\H שלום}, 
% and here's some Japanese: {\J 今日は}.

